\documentclass{article}
\usepackage[utf8]{inputenc}
\usepackage[T1]{fontenc}
\usepackage[margin=2cm]{geometry}
\usepackage{hyperref}
\sloppy
\begin{document}
\section*{Historical source}
Encyclopaedia Britannica scan, source row 2087. Parse 3.5 extracted the text. Super 3.5 VL corrected it from an enlarged crop of each block. It can still contain recognition errors.
Page headers, footers, and page numbers are removed.

\subsection*{BRITAIN (E0179)}
\textit{Continues an entry from an earlier page.}

Reign of France? This was a point which ought not to have been discussed. The object was to treat upon actual circumstances and the real grounds of dispute. Suppose that Bonaparte, desirous to attain peace by any means, should sit down to consider how he could succeed; what does the note allow him to do? He would find that the restoration of the hereditary line of kings was the only case in which a speedy peace would be admitted as possible; in fact, therefore, this restoration was the sine qua non in which immediate negotiation was admissible with ministers. But surely if the ambition of the republic was so formidable, we could not forget this ground of apprehension when we talked of restoring the house of Bourbon. Had we forgotten their proverbial ambition, and was their restoration the remedy for evils arising from such a source? We had now taken up the principle, so much objected to Jacobins, of distinguishing between the people and their government. But what was the conduct of the French? Bonaparte distinctly renounced this principle in the letter to the king, and acknowledged the title and the character of his majesty's government. The note of our ministers was a manifesto to the royalists, and formed for that purpose alone. Lord Holland further stated, that the people at large disapproved of the abrupt rejection of Bonaparte's overtures; he therefore gave his most decided support to the amendment. The Earl of Carnarvon would not consider the answer of our ministers as a refusal to treat for peace, or a declaration of eternal war; it was a call upon the house and the country to pause before they rashly suffered themselves to enter into a negotiation with an unsettled government. He did not expect any extraordinary faith to be manifested by Bonaparte, more than by any other chief or chiefs; but although he would be best pleased with the restoration of monarchy in France, in all times, in monarchies as well as republics, aristocracies, and every other species of government, good faith in treaties was preserved and exemplified only as long as it was the interest of the parties to maintain it. So little integrity had history left on record, that, at the very time they were signed, a secret intention was often indulged to violate them at a particular period. The address, as moved by Lord Grenville, was then carried by a great majority. In the House of Commons Mr Dundas moved a similar address, which gave rise to a similar debate. Mr Dundas said, that the leading feature of the French revolution was a disregard of all treaties, and a contempt for the rights of other powers; and in proof of this assertion, he considered it as necessary merely to recite the names of Spain, Naples, Sardinia, Tuscany, Genoa, Geneva, Modena, Austria, Russia, England, and Egypt, with Denmark and Sweden, though at all times neutral states. Britain had not at this time any reasonable cause to suppose that a change of principles had taken place. The Jacobinical form of government was indeed at an end; but in substance and essence all the qualities of the revolutionary government were in as full force at this moment as they were in the days of Robespierre. Mr Whitbread asserted, that, had it not been for the interference and ambition of the other powers of Europe, the French revolution would have assumed a very different character from that which it now exhibited. He remarked, that other powers had treated neutral states no less unjustly than had been done by the French; Lord Harvey and Lord Hood had ordered the French ministers to be dismissed from Florence; and by threats we had compelled Genoa to dismiss her French inhabitants. He compared Bonaparte with Suwaroff, and the invasion of Egypt by France with that of Poland by Austria, Russia, and Prussia, whose friendship we had frequently courted. Mr Thomas Erskine entered at great length into the question, whether France or Great Britain

had been guilty of the original aggression in the war, and contended that the British government had engaged in it unnecessarily, and persisted in it without necessity. Mr Pitt, on the other hand, affirmed that the French leaders had themselves begun the war, on the principle that it was necessary to consolidate the revolution. With regard to the proposal to negotiate with their present leader, it was impossible to discuss fairly its propriety, without taking into consideration his personal character and conduct. Mr Pitt then expatiated on the conduct of Bonaparte at Campo Formio, in the Milanese, Genoa, Modena, Tuscany, Rome, Venice, Switzerland, and Egypt. His arts of perfidy, he said, were commensurate with the number of treaties; and if we traced the history of the men in this revolution whose conduct had been marked by the most atrocious cruelty, the name of Bonaparte would be found allied to more of them than that of any other within these ten eventful and disastrous years. From these facts the house might judge what reliance might reasonably be placed on this conqueror, and what degree of credit might be given to his professions. It had been observed, indeed, that whatever had been his character, he had now an interest in making and preserving peace. This was a doubtful proposition. That it was his interest to negotiate, would be readily acknowledged, and to negotiate with this country separately, in order to dissolve the whole system of the confederacy on the Continent; to parleyse at once the arms of Russia, of Austria, or of any other country which might look to us for support; and then either to break off his separate treaty, or, if he should have concluded it, to apply the lesson taught in his school of policy in Egypt, and to revive at his pleasure those claims of indemnification which may have been reserved to some happier period. Under all these circumstances of his personal character and his newly acquired power, what security had he for retaining that power but the sword? His hold upon France was the sword, and he had no other. But was the inference to be drawn from these considerations, that we ought in no case to treat with Bonaparte? No; but we ought to wait for the evidence of facts. At present there was nothing from which we could presage a favourable disposition in the French consuls. There was the greatest reason to rely on powerful co-operation from our allies; the strongest indication in the interior of France of a disposition to resist this new tyranny; and every ground to believe, that if we were disappointed of complete success, the continuance of the contest, instead of making our situation comparatively worse, would make it comparatively better. With regard to the negotiation at Lisle in 1797, to which allusions had been made, the Jacobin system of prodigality and bloodshed, by which the efforts of France had been supported, had at that period driven us to exertions which had exhausted the ordinary means of defraying our immense expenditure, and led many who were convinced of the necessity of the war to doubt the possibility of persisting in it. Under this impression we negotiated, not from the sanguine hope that its result would be permanent security, but from the persuasion that the danger arising from peace in these circumstances would be less than the continuance of war with inadequate means. Mr Fox was decidedly of opinion that France, at the commencement of the war, was the defending party. The aggressions of Austria and Prussia could not be denied by any impartial person: nothing could be more decidedly hostile than their proceedings; they scrupled not to declare to France that it was her internal concerns, not her outward actions, which provoked them to confederate against her; they did not pretend to fear her ambition, her conquests, her troubling her neighbours; but they accused her of new-modelling her own government. In all this he


\section*{Source record}
\begin{description}
\item[Scan ID] \texttt{britannica:309f765b1ecc83e0a7f7762e44b4a927a7555443:train:2087}
\item[Image file] \texttt{page-2087-4b834c6b315d.jpg}
\item[Image SHA-256] \texttt{4b834c6b315dd646994b9bac6e843fe2108b703185b904d99893a4a42ef304b1}
\item[Dataset] NationalLibraryOfScotland/encyclopaedia\_britannica\_illustrated at revision \texttt{309f765b1ecc83e0a7f7762e44b4a927a7555443} (CC0-1.0)
\item[Extraction] model \texttt{nvidia/nemotron-3.5-parse}, raw reply \texttt{extract/raw/page-2087-4b834c6b315d.json}
\item[Raw reply SHA-256] \texttt{1296896d5c793a57cb7fa2be42dc32d3435f21bfcf9969f7d7c67cc52b5215ed}
\item[Correction] model \texttt{nvidia/Nemotron-3.5-Super-VL}, block crops, raw replies \texttt{correct/raw/crop\_nothink/page-2087-4b834c6b315d-b*.json}; super\_vl: 2
\item[Running header] BRITAIN. 571
\end{description}
\end{document}
